\section{TASK 4: INFERENCE} \label{sec:task4}

The final stage applies the best validation model to the unlabelled
\texttt{cyberlab.csv} sessions. The selected model is \textbf{UnixCoder}, which
has the highest validation macro F1 among the fully trained and frozen
configurations. The same word truncation used in Task~2 and Task~3 is applied
before inference. Since the inference set is unlabelled, no test metric is
computed; the model output is used only for qualitative and aggregate analysis.
For each Bash word, the predicted tactic is taken from the \textbf{first
sub-token}; if a session is truncated, only words that receive a prediction are
considered.
The inference dataset contains \textbf{174,262 sessions}, spanning from
\textbf{2019-09-01} to \textbf{2019-12-30}.

\subsection{Command-Level Tactic Profiles} \label{subsec:task4_commands}

Table~\ref{tab:task4_commands} reports the predicted tactic distribution for
the required commands. The mapping is not one-to-one. \texttt{grep} is almost
uniquely associated with \textbf{Discovery}, which is coherent with its common
use for inspecting system information. \texttt{cat} is also mainly Discovery,
but it appears as Execution when it belongs to command chains that decode or run
payloads. \texttt{rm} is mostly Discovery or Execution, with smaller Defense
Evasion and Persistence components. \texttt{echo} is the most context-dependent
command: it is predicted as Persistence when it writes credentials or startup
artifacts, as Execution when it feeds payloads to the shell, and as Discovery
when used in benign-looking environment checks.

\begin{table}[H]
    \centering
    {\footnotesize
    \renewcommand{\arraystretch}{1.08}
    \setlength{\tabcolsep}{7pt}
    \begin{tabular}{@{}llrr@{}}
        \toprule
        \textbf{Command} & \textbf{Predicted tactic} & \textbf{Count} &
        \textbf{Freq. (\%)} \\
        \midrule
        \texttt{cat} & Discovery & 724,370 & 84.07 \\
        \texttt{cat} & Execution & 137,283 & 15.93 \\
        \texttt{cat} & Not Malicious Yet & 1 & 0.00 \\
        \texttt{echo} & Persistence & 327,076 & 42.67 \\
        \texttt{echo} & Discovery & 252,994 & 33.01 \\
        \texttt{echo} & Execution & 186,177 & 24.29 \\
        \texttt{echo} & Other & 190 & 0.02 \\
        \texttt{echo} & Impact & 1 & 0.00 \\
        \texttt{grep} & Discovery & 996,075 & 99.97 \\
        \texttt{grep} & Persistence & 302 & 0.03 \\
        \texttt{rm} & Discovery & 248,848 & 72.14 \\
        \texttt{rm} & Execution & 82,502 & 23.92 \\
        \texttt{rm} & Defense Evasion & 8,180 & 2.37 \\
        \texttt{rm} & Persistence & 5,442 & 1.58 \\
        \bottomrule
    \end{tabular}
    \vspace{0.1cm}
    \caption{Predicted MITRE tactic frequencies for selected commands in the
    inference set.}
    \label{tab:task4_commands}
    }
\end{table}

\noindent
Table~\ref{tab:task4_examples} lists one
representative session for each unique (command, predicted tactic) tuple
observed in Table~\ref{tab:task4_commands}. The examples are the first matching
session encountered during inference. Since sessions typically contain multiple
target commands, the same session naturally appears for several tuples.

\vspace{0.1cm}

\begin{table}[H]
    \centering
    \scriptsize
    \renewcommand{\arraystretch}{1.12}
    \begin{tabular}{@{}llp{9.5cm}@{}}
        \toprule
        \textbf{Cmd} & \textbf{Predicted tactic} & \textbf{Example session (truncated)} \\
        \midrule
        \texttt{cat} & Discovery & {\ttfamily enable ; system ; shell ; sh ; cat /proc/mounts; /bin/busybox TIPZU ; cd /dev/shm; cat .s || cp /bin/echo .s; /bin/busybox TIPZU ; tftp; wget; /bin/busybox TIPZU ; dd bs=52 count=1 if=.s || cat .s || while read i; do ech...} \\
        \texttt{cat} & Execution & {\ttfamily cat /proc/cpuinfo | grep name | wc -l ; echo "root:XP3IUReH9hhH"|chpasswd|bash ; echo "321" > /var/tmp/.var03522123 ; rm -rf /var/tmp/.var03522123 ; cat /var/tmp/.var03522123 | head -n 1...} \\
        \texttt{cat} & Not Malicious Yet & {\ttfamily cat /etc/issue...} \\
        \addlinespace
        \texttt{echo} & Persistence & \textit{Same session as} \texttt{cat}$\,\to\,$\textit{Execution} \\
        \texttt{echo} & Discovery & \textit{Same session as} \texttt{cat}$\,\to\,$\textit{Discovery} \\
        \texttt{echo} & Execution & \textit{Same session as} \texttt{cat}$\,\to\,$\textit{Execution} \\
        \texttt{echo} & Other & {\ttfamily cd /tmp || /var/tmp || /dev/shm; echo "ZXZhbCB1bnBhY2sgdT0+cXtfIkZVWSgiMVA8Rl1...} \\ %DOTctUztSYF0oIj1SPFdFTjhSPFsi...} \\
        \texttt{echo} & Impact & {\ttfamily echo "Z""IGAZAGA148""8" ; cat /dev/stdin | sh ; \#!/usr/bin/sh ; \#timeout=30 ; killSTDIN() \{ ; processID=\$(ps -u | grep -v 'grep' | grep cat...} \\
        \addlinespace
        \texttt{grep} & Discovery & \textit{Same session as} \texttt{cat}$\,\to\,$\textit{Execution} \\
        \texttt{grep} & Persistence & {\ttfamily scp -t \textasciitilde{}/6vqzhh1w2sus12ms9e1oxuqryj ; LC\_ALL=C \textasciitilde{}/6vqzhh1w2sus12ms9e1oxuqryj ; LC\_ALL=C rm -f \textasciitilde{}/6vqzhh1w2sus12ms9e1oxuqryj ; LC\_ALL=C chattr -i -a \textasciitilde{}/.dhpcd...} \\
        \addlinespace
        \texttt{rm} & Defense Evasion & \textit{Same session as} \texttt{cat}$\,\to\,$\textit{Discovery} \\
        \texttt{rm} & Discovery & \textit{Same session as} \texttt{cat}$\,\to\,$\textit{Execution} \\
        \texttt{rm} & Execution & \textit{Same session as} \texttt{cat}$\,\to\,$\textit{Execution} \\
        \texttt{rm} & Persistence & {\ttfamily cat /proc/cpuinfo | grep name | wc -l ; echo "root:9R8X8aKPE6l6"|chpasswd|bash ; echo "321" > /var/tmp/.var03522123 ; rm -rf /var/tmp/.var03522123...} \\
        \bottomrule
    \end{tabular}
    \vspace{0.1cm}
    \caption{One example session per (command, predicted tactic) tuple.
    Sessions containing multiple target commands appear for more than one tuple.}
    \label{tab:task4_examples}
\end{table}

\noindent
The examples confirm that the same Bash word receives different tactic labels
depending on the surrounding command context. For instance,
\texttt{cat /proc/cpuinfo | grep name} follows a Discovery pattern, while
\texttt{echo <payload> | bash} maps naturally to Execution.
\texttt{echo root:<password> | chpasswd | bash} supports Persistence because it
modifies access credentials, and \texttt{rm -rf} inside cleanup chains is
labelled as Defense Evasion. The rare tuples (e.g.\ \texttt{cat}$\to$Not
Malicious Yet, \texttt{echo}$\to$Impact) are isolated contextual decisions
inside long chains rather than systematic associations.

\subsection{Fingerprint Analysis} \label{subsec:task4_fingerprints}

Each session is converted into a fingerprint, defined as the ordered sequence
of predicted MITRE tactics. The inference set produces \textbf{1,905 unique
fingerprints}; \textbf{901} of them occur only once, indicating a long tail of
rare behaviours. Fingerprints are sorted by date of birth, assigned increasing
IDs, and counted per day.

Before looking at the temporal structure, the frequency distribution already
shows that most fingerprints are rare, whereas a small number of recurrent
patterns account for a large portion of the traffic. This is expected in
honeypot data: many sessions are minor variants of automated scripts, while
occasional payloads or command permutations create one-off fingerprints. The
total number of unique fingerprints is sensitive to per-token prediction
consistency: a model that tags structurally similar bot-generated sessions more
uniformly (as UnixCoder does, given its higher macro F1) collapses many
near-duplicate sessions into a smaller set of recurrent templates, while a
noisier model would fragment them into more near-duplicate fingerprints.

\vspace{0.3cm}

\begin{figure}[!htbp]
    \centering
    \includegraphics[width=\linewidth]{\imagespath Task4/fingerprint_frequency_distribution.png}
    \caption{Distribution of fingerprint frequencies and most common inferred
    fingerprints.}
    \label{fig:task4_fingerprint_frequency}
\end{figure}

\vspace{0.2cm}

\noindent
Figure~\ref{fig:task4_fingerprint_frequency} makes this imbalance explicit. The
left panel highlights the long tail of low-frequency fingerprints, while the
right panel shows that the most common patterns are dominated by repeated
Discovery-heavy sequences. Therefore, a fingerprinting view is useful because it
compresses a large inference set into a smaller number of behavioural templates.

Figure~\ref{fig:task4_fingerprint_scatter} shows the required time-vs-ID view.
Older fingerprints appear near the bottom of the plot. No fingerprint is
literally present during the entire \textbf{121-day} observation window, but
some early fingerprints are highly persistent: IDs \textbf{0}, \textbf{11}, and
\textbf{3} are active for 99, 98, and 96 days respectively. The most common
fingerprint appears in \textbf{19,247 sessions}.

\begin{figure}[!htbp]
    \centering
    \includegraphics[width=0.96\linewidth]{\imagespath Task4/fingerprint_scatter_plot.png}
    \caption{Fingerprint IDs over time. Point size and colour encode the number
    of sessions assigned to a fingerprint on a given day.}
    \label{fig:task4_fingerprint_scatter}
\end{figure}

The scatter plot also exposes concentrated bursts that are compatible with
automated campaigns. The adaptive detector finds \textbf{6 campaign spikes}.
The strongest burst is fingerprint \textbf{1594}, with \textbf{1,736 sessions}
on \textbf{2019-12-09} and related spikes on 2019-12-08 and 2019-12-10.
Fingerprint \textbf{293} shows a similar December concentration, exceeding
900 sessions per day in the same period. These high-volume fingerprints include
repeated Discovery chains such as CPU inspection followed by credential or
payload-related commands, suggesting repeated scripted activity rather than
isolated manual interaction.

\begin{figure}[!htbp]
    \centering
    \includegraphics[width=0.95\linewidth]{\imagespath Task4/daily_volume_and_novelty.png}
    \vspace{0.2cm}
    \caption{Daily session volume and novelty dynamics for the inferred
    fingerprints.}
    \label{fig:task4_daily_volume}
\end{figure}

Figure~\ref{fig:task4_daily_volume} provides a complementary aggregate view. The
daily volume is not uniform: activity grows after mid-October and reaches its
largest peak in early December, in agreement with the fingerprint spikes
identified above. The novelty curve is higher at the beginning of the
collection, when many fingerprints are observed for the first time, and becomes
lower once repeated behaviours dominate the traffic. This suggests two
different regimes: an initial phase with many newly observed patterns, followed
by a period in which recurrent automated behaviours account for a large share
of the sessions.

\begin{figure}[H]
    \centering
    \vspace{1.5cm}
    \includegraphics[width=0.95\linewidth]{\imagespath Task4/daily_tactic_volume_and_composition.png}
    \vspace{0.2cm}
    \caption{Daily MITRE tactic volume and normalized tactic composition in the
    inference set.}
    \label{fig:task4_daily_tactics}
\end{figure}

\noindent
The tactic-level view in Figure~\ref{fig:task4_daily_tactics} confirms that the
traffic volume is largely driven by Discovery-labelled words, but the relative
composition is not perfectly constant. Execution and Persistence occupy visible
shares during the central part of the observation window, while the strongest
December peak is mostly a scale effect: many more sessions are observed, but
the dominant tactical structure remains similar to the repeated fingerprints
discussed above. This is useful from an analyst perspective because it
distinguishes between a completely new behavioural mix and a large-scale replay
of already observed attack patterns.
